\documentclass[11pt,a4paper]{article}
\usepackage[margin=25mm]{geometry}
\usepackage{booktabs,longtable,tabularx,array}
\usepackage{amsmath}
\usepackage{hyperref}
\usepackage{xcolor}
\usepackage{enumitem}
\usepackage{microtype}
\title{HypatiaX Core-15 LLM Ablation\\Performance Evaluation Report}
\author{LLM-HypatiaX-REPRO}
\date{26 November 2026}
\begin{document}
\maketitle
\begin{abstract}
This report evaluates the merged output of GitHub Actions run
\href{https://github.com/sednabcn/LLM-HypatiaX-REPRO/actions/runs/45869191966}{45869191966}
(commit \texttt{b460be2475ad9527bb7a56c9aaa11e740d26dea6}).
The experiment compares PySR-only, HypatiaX with its LLM component, and
HypatiaX with the LLM disabled on 15 equations. The workflow completed
successfully, but execution success is not evidence of scientific superiority.
Median held-out test accuracy is essentially tied. HypatiaX has a lower
observed median wall-clock time, but timing is heterogeneous and the arms
were not run under identical computational budgets. Far-extrapolation
results are mixed, with two non-finite/missing HypatiaX cases excluded from
the paired extrapolation analysis. The reported paired test does not establish
a statistically significant HypatiaX advantage ($p=0.0874$).
\end{abstract}
\section{Scope and provenance}
The analyzed artifact is \texttt{results-exp1\_ablation-merged}, produced by
run 45869191966. The merged JSON contains 15 equation records and the
artifact contains 94 files, including checkpoints, shard outputs, logs,
statistical summaries, provenance and LaTeX tables. The run used seed 42,
30 populations and engine version 5.4 according to the merged metadata.
Five shards are workload partitions, not five independent seed replications.
\section{Experimental conditions}
\begin{itemize}[leftmargin=*]
\item PySR-only: symbolic-regression baseline.
\item HypatiaX: hybrid system with LLM-assisted candidate generation/warm-start.
\item HypatiaX-noLLM: hybrid system with the LLM disabled.
\end{itemize}
The merged records report 15 completed records in each arm. LLM was called
for all 15 HypatiaX records. The primary far-extrapolation paired analysis
contains 13 equations; Arrhenius and Impermanent Loss are excluded because
the HypatiaX far-extrapolation result is non-finite/unavailable.
\section{Aggregate results}
\begin{table}[ht]
\centering
\caption{Descriptive aggregate metrics across Core-15.}
\begin{tabular}{lrrr}
\toprule
Metric & PySR-only & HypatiaX & HypatiaX-noLLM\\
\midrule
Median test $R^2$ (n=15) & 0.99724 & 0.99722 & 0.99724\\
Median test RMSE (n=15) & 0.08365 & 0.08374 & 0.08374\\
Median far-OOD $R^2$ & 0.99994 (n=15) & 0.99971 (n=13) & 0.99957 (n=13)\\
Mean far-OOD $R^2$ & -7.2545 (n=15) & 0.3387 (n=13) & 0.2254 (n=13)\\
Median observed runtime (s) & 1101.30 & 362.49 & 385.71\\
Median expression complexity & 7 & 33 & 33\\
\bottomrule
\end{tabular}
\end{table}
The far-OOD medians must not be compared without noting their different
denominators. Excluding failed/non-finite cases can make a median look
reassuring while hiding important reliability failures. The mean far-OOD
values are highly sensitive to severe negative $R^2$ outliers.
Negative $R^2$ means the model is worse than the constant test-mean predictor.
The observed HypatiaX median runtime is about 3.04 times lower than PySR-only
(1101.30/362.49). This is descriptive, not a controlled speedup claim:
PySR-only has a 1400-second budget and one attempt, while HypatiaX has a
3600-second budget and up to three attempts in the recorded configuration.
Some HypatiaX cases complete in only a few seconds, while Gravitational Force
takes about 1937 seconds. A matched-budget, repeated comparison is needed.
\section{Far-extrapolation paired test}
The saved primary analysis reports $n=13$, Wilcoxon signed-rank
$W=17$, one-sided $p_{H>P}=0.0874$, and paired rank-biserial effect size
$r=+0.48$. The predeclared 0.05 threshold is not crossed. The Mann--Whitney
sensitivity analysis reports $U=89$, one-sided $p=0.4183$.
The observed direction favours HypatiaX on three equations, with ten ties
under the analysis' tie tolerance and no PySR-only wins. This is suggestive
but not statistically significant evidence.
\section{Equation-level findings}
\begin{longtable}{p{35mm}p{30mm}p{85mm}}
\toprule
Equation & Domain & Interpretation\\
\midrule
Arrhenius & Chemistry & HypatiaX far-OOD result is unavailable/non-finite; near-OOD performance is poor. The saved analysis flags an LLM-prior/search-space failure mode.\\
Henderson--Hasselbalch & Chemistry & Large improvement relative to PySR-only in far-OOD $R^2$ (approximately $-34.47$ to $0.999$).\\
Rate Law & Chemistry & Both methods perform strongly; PySR-only is marginally higher in far-OOD.\\
Allometric Scaling & Biology & Both methods reach near-perfect far-OOD scores.\\
Michaelis--Menten & Biology & Both methods degrade under far extrapolation; HypatiaX improves from about 0.084 to 0.184 but remains weak.\\
Logistic Growth & Biology & Both methods perform very strongly.\\
Kinetic Energy & Physics & Both methods perform very strongly.\\
Gravitational Force & Physics & Both methods fail badly in far extrapolation; HypatiaX is less negative but still poor.\\
Ideal Gas Law & Physics & Both methods perform very strongly.\\
Impermanent Loss & DeFi AMM & HypatiaX far-OOD result unavailable/non-finite; near/medium extrapolation also exposes instability.\\
Price Impact & DeFi AMM & Both methods perform very strongly.\\
Constant Product & DeFi AMM & Both methods perform very strongly.\\
Value at Risk & DeFi Risk & Both methods perform very strongly.\\
Liquidation Price & DeFi Risk & Both methods perform very strongly.\\
Portfolio Std Dev & DeFi Risk & Both methods perform very strongly.\\
\bottomrule
\end{longtable}
\section{LLM contribution}
The saved LLM-contribution analysis reports 15 LLM calls, with HypatiaX and
HypatiaX-noLLM far-OOD paired analysis on 13 equations. The Holm-adjusted
one-sided Wilcoxon $p$-value is 0.375 for each of the three reported
comparisons. There is no statistically significant evidence in this run
that adding the LLM improves far-OOD performance over the no-LLM variant.
The final selected expression is identical between HypatiaX and no-LLM in
8 of 15 equation records, indicating that LLM invocation does not necessarily
change the final selected model.
\section{Validity threats and limitations}
\begin{enumerate}[leftmargin=*]
\item Only one random seed (42) is represented; five shards are not independent replicates.
\item Two HypatiaX far-OOD results are unavailable/non-finite and must be counted as failures in reliability reporting, not silently dropped.
\item The compute budgets and allowed attempt counts differ between baseline and hybrid arms.
\item Expression complexity is higher for HypatiaX (median 33 versus 7), which may affect interpretability and overfitting risk.
\item Very high scores on several equations and very poor scores on a few require per-equation reporting; averages alone are misleading.
\item This report evaluates the single \texttt{exp1\_ablation} run, not all repository experiments.
\end{enumerate}
\section{Conclusions}
The current evidence supports three restrained conclusions:
\begin{itemize}[leftmargin=*]
\item HypatiaX achieves held-out predictive accuracy comparable to PySR-only on this run.
\item Its observed median runtime is lower, but a controlled equal-budget experiment is required before claiming a speed advantage.
\item The run does not establish statistically significant far-extrapolation superiority or a significant incremental benefit from the LLM.
\end{itemize}
The main research risk is robustness: a small number of severe extrapolation failures and non-finite outcomes coexist with excellent results on many other equations.
\section{Recommended evaluation programme}
\begin{enumerate}[leftmargin=*]
\item \textbf{Multi-seed replication.} Repeat all three arms with seeds
42, 99, 123, 777 and 2024. Pair methods by equation and seed. Report mean,
median, standard deviation, confidence intervals, win/tie/loss counts and
failure rates. Treat seed as a repeated-measures factor; do not count shards
as replicates.
\item \textbf{Matched compute budgets.} Define one common wall-clock or CPU
budget, attempt policy, hardware/software environment and stopping rule for
all arms. Record actual CPU/GPU time, wall time, number of attempts and
resource usage. Include both equal-budget and native-budget sensitivity runs.
\item \textbf{Complete failure-aware reporting.} Preserve all 15 equations
in tables. Report non-finite/failed predictions explicitly; provide success
rates and worst-case/quantile metrics in addition to conditional metrics
among successful cases. Keep near, medium and far OOD regimes separate.
\item \textbf{LLM ablation.} Compare HypatiaX with LLM against the identical
pipeline with LLM disabled, holding seeds, budgets, data splits, search
settings and retry policy fixed. Log LLM calls, latency, token/cost, proposed
formula, acceptance/rejection reason and whether the selected expression
changed. Use paired per-equation/per-seed analysis with Holm correction.
\item \textbf{Complexity and scientific validity.} Add symbolic complexity,
dimensional consistency, domain constraints, singularity checks and
equivalence-to-ground-truth checks where possible. Report a Pareto frontier
of accuracy, extrapolation, runtime and complexity rather than accuracy alone.
Have domain-specific checks pre-registered before examining outcomes.
\end{enumerate}
\section{Practical implementation sequence}
First create a benchmark manifest fixing the 15 equations, train/test/OOD
splits, five seeds and arm definitions. Next add a multi-seed matrix or
orchestrator that dispatches one run per (seed, arm), with explicit seed
pass-through and unique output directories. Then add a common budget
configuration and persist actual resource accounting. Extend aggregation
to retain failed records and produce a complete equation-by-seed-by-arm
table. Finally add paired statistical analysis and complexity/scientific
validity diagnostics. Run a small smoke subset before the full benchmark;
the full multi-seed campaign should be launched only after configuration
review and explicit approval.
\section{Reproducibility references}
\begin{itemize}[leftmargin=*]
\item Workflow run: \url{https://github.com/sednabcn/LLM-HypatiaX-REPRO/actions/runs/45869191966}
\item Commit: \texttt{b460be2475ad9527bb7a56c9aaa11e740d26dea6}
\item Primary source files in the merged artifact: \texttt{exp1\_ablation\_results.json},
\texttt{exp1\_rf01\_mannwhitney.json}, \texttt{exp1\_llm\_contribution.json},
\texttt{exp1\_instability\_stats.json}, and \texttt{wall\_clock\_flags.json}.
\end{itemize}
\end{document}